\section{问题重述}

\subsection{问题背景}
作为古代中西方文明交流的核心通道，丝绸之路承载了大量物质与技术的跨区域传播，玻璃器物作为手工业技术与贸易往来的直接载体，是还原该历史进程的关键实物证据。传统玻璃烧制多以石英砂为基础原料，辅以助熔剂降低熔点，并添加稳定剂提升化学稳定性。受产地原料禀赋与制作工艺差异影响，外观相近的玻璃文物往往呈现出截然不同的化学组成特征。本次研究对象涵盖高钾玻璃与铅钡玻璃两大体系，二者因助熔剂体系迥异，导致氧化钾、氧化铅、氧化钡等关键组分含量存在显著分异。

玻璃文物经长期地下埋藏，普遍存在不同程度的表面风化现象；埋藏环境中的元素与文物本体发生离子交换，导致表层化学成分占比发生偏移，为玻璃材质的准确判别带来干扰。题目附件提供了58件已标定类别的玻璃文物基础信息、69个采样点位的14项化学成分占比数据，以及编号为A1至A8的8件未知类别文物的检测结果。数据有效性判定标准为：各组分比例总和处于85\%～105\%区间内视为有效数据，空白检测项表示该成分含量低于检出限。

\subsection{问题提出}
围绕上述检测数据，本文需依次解决四项核心问题：
\begin{enumerate}
  \item 探究表面风化程度与玻璃类别、纹饰特征及颜色属性之间的关联规律；分类型对比风化前后的化学成分演化特征，并对风化样品的原始组分进行反演推算。
  \item 提炼高钾玻璃与铅钡玻璃的分类判别准则；分别选取适宜的化学成分指标开展亚类划分工作，并对分类结果的合理性与敏感性进行检验。
  \item 对8件未知类别文物样品进行玻璃类型判别，并评估分类结果的稳定程度。
  \item 分类型剖析两类玻璃体系内部各化学成分的相关性特征，并对比二者在成分关联模式上的差异。
\end{enumerate}

\subsection{研究目标}
本文研究面向存在成分比例约束、风化干扰及多点位采样特征的古代玻璃检测数据，旨在构建一套系统化的成分分析与类别判别方法，而非对文物进行价值优劣评判；具体涵盖统计推断、风化前成分反演、有监督分类、无监督聚类及关联结构对比等核心任务。